\documentclass[a4paper,12pt]{article}
\usepackage{color}
\usepackage{graphicx, gensymb}
\usepackage{amsfonts, cancel, amssymb}
\usepackage{amsmath, amsthm, array, esvect, esint}
\usepackage{typearea}
\usepackage[a4paper,left=2cm,right=2cm,top=2cm,bottom=3cm,bindingoffset=5mm]{geometry}
\newcommand\numberthis{\addtocounter{equation}{1}\tag{\theequation}}
\newcommand{\bra}{}
\newcommand{\kom}{}
\newcommand{\brak}{}
\newcommand{\braket}{}
\newcommand{\intx}{}
\newcommand{\intr}{}
\newcommand{\kla}{}
\newcommand{\klam}{}
\newcommand{\klammer}{}
\newcommand{\oper}{}
\newcommand{\tecxt}{}
\newcommand{\Bra}{}
\newcommand{\Ket}{}

\begin{document}

\title{11 Eichinvarianz und Higgs-Boson}
\author{Joachim Schwardt}
\date{\today}
\maketitle

\section*{Aufgabe 11.1: Lokale Eichinvarianz}
Die Dirac-Gleichung ist $(i\gamma^\mu\partial_\mu - m)\psi=0$. Sie ist trivialerweise invariant unter einer globalen Transformation $\psi\rightarrow\psi':=e^{i\theta}\psi$ mit $\theta=\tecxt\text{const}$:
\begin{align*}
(i\gamma^\mu\partial_\mu - m)\psi'&=(i\gamma^\mu\partial_\mu - m)e^{i\theta}\psi=e^{i\theta}(i\gamma^\mu\partial_\mu - m)\psi=0
\end{align*}
Für eine lokale Transformation $\psi\rightarrow\psi':=e^{iq\chi(x)}\psi$ gilt das nicht mehr:
\begin{align*}
(i\gamma^\mu\partial_\mu - m)\psi'&=e^{iq\chi}(i\gamma^\mu\partial_\mu - m)\psi + \psi (i\gamma^\mu\partial_\mu ) e^{iq\chi}=-q\psi'\gamma^\mu\partial_\mu\chi\neq 0
\end{align*}
Führt man noch die minimale Substitution $\partial_\mu\rightarrow D_\mu :=\partial_\mu +iqA_\mu$, so kann die Invarianz wiederhergestellt werden. Dazu muss noch angenommen werden, dass für $A_\mu$ die Transformation $A_\mu\rightarrow A_\mu'=A_\mu - \partial_\mu\chi$ gilt:
\begin{align*}
D_\mu'\psi'&=(\partial_\mu +iqA_\mu')e^{iq\chi}\psi=\\
&=e^{iq\chi}\partial_\mu\psi + \psi\partial_\mu e^{iq\chi} + iq(A_\mu - \partial_\mu\chi)e^{iq\chi}\psi=\\
&=e^{iq\chi}\kla\left( {\partial_\mu +iqA_\mu } \right)\psi +\psi e^{iq\chi}\partial_\mu (iq\chi ) - iqe^{iq\chi}\psi\partial_\mu\chi=\\
&=e^{iq\chi}D_\mu\psi
\end{align*}
Daraus folgt direkt die Invarianz der Dirac-Gleichung:
\begin{align*}
(i\gamma^\mu D_\mu' - m)\psi'&=e^{iq\chi}(i\gamma^\mu D_\mu - m)\psi=0
\end{align*}
\section*{Aufgabe 11.2: Elektroschwache Parameter}
Es gelten folgende Formeln, wobei $\alpha_\tecxt\text{QED}\simeq\frac{1}{137}$:
\begin{align*}
m_W&=\frac{1}{2}vg_W=80.4\,\tecxt\text{GeV}\numberthis\label{mW}\\
m_Z&=\frac{1}{2}v\sqrt{g_W^2+g'^2}=91.2\,\tecxt\text{GeV}\numberthis\label{mZ}\\
\cos\theta_W&=\frac{g_W}{\sqrt{g_W^2+g'^2}}\numberthis\label{cos thetaW}\\
e&=g'\cos\theta_W=\sqrt{4\pi\alpha_\tecxt\text{QED}}\numberthis\label{alpha}
\end{align*}
Daraus lassen sich die elektroschwachen Parameter bestimmen:
\begin{align*}
\frac{(\ref{mW})}{(\ref{mZ})}&&:&&\frac{m_W}{m_Z}&=\cos\theta_W&\Rightarrow &&\theta_W&=\arccos\frac{m_W}{m_Z}=28.17\,\degree\\
(\ref{cos thetaW})\,\tecxt\text{in\,}(\ref{alpha})&&:&&g'&=\frac{\sqrt{4\pi\alpha_\tecxt\text{QED}}}{\cos\theta_W}&\Rightarrow &&g'&=\sqrt{4\pi\alpha_\tecxt\text{QED}}\frac{m_Z}{m_W}=0.344\\
(\ref{mZ})^2-(\ref{mW})^2&&:&&\frac{1}{4}v^2g'^2&=m_Z^2-m_W^2&\Rightarrow &&v&=\frac{2}{g'}\sqrt{m_Z^2-m_W^2}=250.6\,\tecxt\text{GeV}\\
(\ref{mW})&&:&&g_W&=\frac{2m_W}{v}&\Rightarrow &&g_W&=\frac{g'm_W}{\sqrt{m_Z^2-m_W^2}}=0.642
\end{align*}

\end{document}